\section{End-Term Comprehensive Examination --- Set B (70 Marks)}
\label{sec:endterm_set_b}

\LPUPaperHeader{End-Term Examination}{B}{3 Hours}{70 Marks}{Comprehensive All Units CO1 to CO6}

\begin{center}
\sffamily\textbf{\color{AlertRed}Structure: Part A (30 Marks Objective MCQs) + Part B (40 Marks Descriptive Analytical)}
\end{center}

\subsection*{Part A: Objective Assessment (30 Marks)}
\textit{Instructions: Answer all 30 questions. Each carries 1 mark. Negative marking: $-0.25$ marks.}

\begin{enumerate}[label=\textbf{\arabic*.}]
    \item {Which is a possible effect of excessive exposure to ionizing radiation?}
    \begin{enumerate*}[label=(\alph*)]
        \item {Faster water purification} \quad
        \item {Damage to body cells} \quad
        \item {Improved hearing ability} \quad
        \item {Increased soil fertility}
    \end{enumerate*}
    \item {Which substance is commonly classified as an emerging pollutant?}
    \begin{enumerate*}[label=(\alph*)]
        \item {Pharmaceutical residue} \quad
        \item {Clean river sand} \quad
        \item {Atmospheric oxygen} \quad
        \item {Natural water vapor}
    \end{enumerate*}
    \item {What is a major environmental effect of fireworks?}
    \begin{enumerate*}[label=(\alph*)]
        \item {Increased forest cover} \quad
        \item {Improved soil quality} \quad
        \item {Air and noise pollution} \quad
        \item {Groundwater recharge}
    \end{enumerate*}
    \item {Which gas is the main greenhouse gas released by burning fossil fuels?}
    \begin{enumerate*}[label=(\alph*)]
        \item {Hydrogen} \quad
        \item {Helium} \quad
        \item {Nitrogen} \quad
        \item {Carbon dioxide}
    \end{enumerate*}
    \item {Which is a likely effect of climate change?}
    \begin{enumerate*}[label=(\alph*)]
        \item {Complete removal of storms} \quad
        \item {Equal rainfall in all regions} \quad
        \item {Permanent weather stability} \quad
        \item {More frequent heat waves}
    \end{enumerate*}
    \item {Which group of chemicals is mainly associated with ozone layer depletion?}
    \begin{enumerate*}[label=(\alph*)]
        \item {Amino acids} \quad
        \item {Chlorofluorocarbons} \quad
        \item {Natural sugars} \quad
        \item {Carbohydrates}
    \end{enumerate*}
    \item {Which gases are mainly responsible for the formation of acid rain?}
    \begin{enumerate*}[label=(\alph*)]
        \item {Methane and argon} \quad
        \item {Oxygen and hydrogen} \quad
        \item {Sulfur dioxide and nitrogen oxides} \quad
        \item {Helium and neon}
    \end{enumerate*}
    \item {The Bhopal gas tragedy was caused by the accidental release of which toxic gas?}
    \begin{enumerate*}[label=(\alph*)]
        \item {Carbon dioxide} \quad
        \item {Nitrogen gas} \quad
        \item {Methyl isocyanate} \quad
        \item {Water vapor}
    \end{enumerate*}
    \item {What do the three Rs of solid waste management represent?}
    \begin{enumerate*}[label=(\alph*)]
        \item {Remove, replace, and refill} \quad
        \item {Recover, repair, and return} \quad
        \item {Record, report, and review} \quad
        \item {Reduce, reuse, and recycle}
    \end{enumerate*}
    \item {Which measure is suitable for controlling industrial wastewater pollution?}
    \begin{enumerate*}[label=(\alph*)]
        \item {Diluting waste with rainwater} \quad
        \item {Using an effluent treatment plant} \quad
        \item {Releasing waste into nearby drains} \quad
        \item {Storing waste beside a river}
    \end{enumerate*}
    \item {What is the main hazardous release associated with a nuclear disaster?}
    \begin{enumerate*}[label=(\alph*)]
        \item {Organic compost} \quad
        \item {Volcanic magma} \quad
        \item {Rainwater runoff} \quad
        \item {Radioactive material}
    \end{enumerate*}
    \item {Which incident is an example of a chemical disaster?}
    \begin{enumerate*}[label=(\alph*)]
        \item {Collapse of a snow slope} \quad
        \item {Leak of toxic gas} \quad
        \item {Shaking of the ground} \quad
        \item {Overflow of a river}
    \end{enumerate*}
    \item {Which event is considered a biological disaster?}
    \begin{enumerate*}[label=(\alph*)]
        \item {Widespread infectious disease outbreak} \quad
        \item {Powerful coastal cyclone} \quad
        \item {Major railway collision} \quad
        \item {Large offshore oil spill}
    \end{enumerate*}
    \item {Which event is a transport accident?}
    \begin{enumerate*}[label=(\alph*)]
        \item {Volcano releasing ash} \quad
        \item {Drought damaging crops} \quad
        \item {River flooding its banks} \quad
        \item {Passenger train derailment}
    \end{enumerate*}
    \item {What is the primary purpose of a national disaster management framework?}
    \begin{enumerate*}[label=(\alph*)]
        \item {Manage international sports events} \quad
        \item {Regulate daily weather changes} \quad
        \item {Coordinate disaster risk and response} \quad
        \item {Set prices for farm products}
    \end{enumerate*}
    \item {In India, which body lays down national policies and guidelines for disaster management?}
    \begin{enumerate*}[label=(\alph*)]
        \item {National Highways Authority} \quad
        \item {University Grants Commission} \quad
        \item {National Disaster Management Authority} \quad
        \item {Election Commission of India}
    \end{enumerate*}
    \item {Which task is commonly performed by governmental agencies during a disaster?}
    \begin{enumerate*}[label=(\alph*)]
        \item {Promoting holiday travel} \quad
        \item {Conducting entertainment programs} \quad
        \item {Advertising commercial products} \quad
        \item {Organizing evacuation and rescue}
    \end{enumerate*}
    \item {How do non-governmental organizations commonly help disaster-affected people?}
    \begin{enumerate*}[label=(\alph*)]
        \item {Operating all military installations} \quad
        \item {Distributing food and essential supplies} \quad
        \item {Changing official weather records} \quad
        \item {Issuing national currency and coins}
    \end{enumerate*}
    \item {Why are community-based organizations valuable during disaster response?}
    \begin{enumerate*}[label=(\alph*)]
        \item {They predict the exact time of disasters} \quad
        \item {They understand local needs and resources} \quad
        \item {They replace all professional rescue teams} \quad
        \item {They control every national emergency agency}
    \end{enumerate*}
    \item {What is an important role of the media during a disaster?}
    \begin{enumerate*}[label=(\alph*)]
        \item {Blocking communication between agencies} \quad
        \item {Hiding official evacuation information} \quad
        \item {Sharing verified warnings and instructions} \quad
        \item {Spreading unconfirmed alarming rumors}
    \end{enumerate*}
    \item {Which Indian law was enacted specifically to protect wild animals, birds, and plants?}
    \begin{enumerate*}[label=(\alph*)]
        \item {Forest Act, 1927} \quad
        \item {Water Act, 1974} \quad
        \item {Wildlife Protection Act, 1972} \quad
        \item {Environment Protection Act, 1986}
    \end{enumerate*}
    \item {The Water (Prevention and Control of Pollution) Act was enacted in which year?}
    \begin{enumerate*}[label=(\alph*)]
        \item {1986} \quad
        \item {1972} \quad
        \item {1974} \quad
        \item {1981}
    \end{enumerate*}
    \item {Which law provides a broad framework for environmental protection in India?}
    \begin{enumerate*}[label=(\alph*)]
        \item {Indian Forest Act, 1927} \quad
        \item {Wildlife Protection Act, 1972} \quad
        \item {Water Act, 1974} \quad
        \item {Environment Protection Act, 1986}
    \end{enumerate*}
    \item {Who led the Bishnoi sacrifice to protect Khejri trees at Khejarli?}
    \begin{enumerate*}[label=(\alph*)]
        \item {Medha Patkar} \quad
        \item {Amrita Devi} \quad
        \item {Sugathakumari} \quad
        \item {Gaura Devi}
    \end{enumerate*}
    \item {What action became the main symbol of the Chipko movement?}
    \begin{enumerate*}[label=(\alph*)]
        \item {Cleaning rivers near cities} \quad
        \item {Hugging trees to prevent felling} \quad
        \item {Planting crops around dams} \quad
        \item {Relocating wildlife from forests}
    \end{enumerate*}
    \item {In which Indian state did the Appiko movement begin?}
    \begin{enumerate*}[label=(\alph*)]
        \item {Karnataka} \quad
        \item {Kerala} \quad
        \item {Odisha} \quad
        \item {Gujarat}
    \end{enumerate*}
    \item {What was a major concern of the Jungle Bachao Andolan?}
    \begin{enumerate*}[label=(\alph*)]
        \item {Building airports within major cities} \quad
        \item {Replacing natural forests with teak plantations} \quad
        \item {Increasing tourism in desert regions} \quad
        \item {Expanding fishing in coastal waters}
    \end{enumerate*}
    \item {The Silent Valley movement opposed a proposed project in which state?}
    \begin{enumerate*}[label=(\alph*)]
        \item {Bihar} \quad
        \item {Kerala} \quad
        \item {Rajasthan} \quad
        \item {Punjab}
    \end{enumerate*}
    \item {Which issue was central to the Narmada Bachao Andolan?}
    \begin{enumerate*}[label=(\alph*)]
        \item {Protection of monuments from air pollution} \quad
        \item {Control of traffic in crowded cities} \quad
        \item {Impacts of large dams on people and nature} \quad
        \item {Expansion of industries in dry regions}
    \end{enumerate*}
    \item {The Tehri Dam was constructed on which river?}
    \begin{enumerate*}[label=(\alph*)]
        \item {Mahanadi} \quad
        \item {Bhagirathi} \quad
        \item {Godavari} \quad
        \item {Yamuna}
    \end{enumerate*}
\end{enumerate}

\subsection*{Part B: Descriptive Analytical Section (40 Marks)}
\textit{Instructions: Answer any four out of the following five 10-mark questions.}

\questionitem{1}{
    (a) State the Ehrlich-Holdren IPAT equation. A city of population $P = 2 \times 10^6$ with per-capita affluence $A = \$5{,}000/\text{yr}$ has technology intensity $T = 0.4\text{ kg } CO_2/\$$. Compute current impact. If population grows by $20\%$ and affluence by $50\%$, find target $T$ for a $20\%$ emissions reduction. [5 Marks]\\
    (b) Formulate the Ghyben-Herzberg coastal aquifer relation. Explain why over-pumping freshwater induces severe saltwater intrusion. [5 Marks]
}{10}{CO1, CO2}{L3 Apply}

\questionitem{2}{
    (a) State Lindeman's 10\% law. Given incident solar radiation of $5.0 \times 10^7\text{ kcal/yr}$ on a forest where GPP is $2.0\%$ and respiration is $40\%$ of GPP, calculate NPP and energy at the tertiary consumer level. [5 Marks]\\
    (b) Contrast the Clementsian stages of primary versus secondary ecological succession. [5 Marks]
}{10}{CO2}{L4 Analyze}

\questionitem{3}{
    (a) Define Norman Myers criteria for Biodiversity Hotspots. Detail the geographic and ecological features of the Western Ghats and Himalayan hotspots. [5 Marks]\\
    (b) An ecological survey records 4 species with abundances $n_1 = 40, n_2 = 30, n_3 = 20, n_4 = 10$ ($N = 100$). Calculate Shannon-Wiener index $H'$ and Simpson's index $1 - D$. [5 Marks]
}{10}{CO3}{L3 Apply}

\questionitem{4}{
    (a) Using Streeter-Phelps oxygen sag formulation, define critical time $t_c$ and critical deficit $D_c$. Why is dissolved oxygen critical for freshwater fisheries? [5 Marks]\\
    (b) Outline the industrial mechanisms of Electrostatic Precipitators (ESPs) and Venturi scrubbers for particulate air pollution control. [5 Marks]
}{10}{CO4}{L4 Analyze}

\questionitem{5}{
    (a) Formulate the Disaster Risk Equation. Detail the 4 Priorities for Action under the Sendai Framework (2015--2030). [5 Marks]\\
    (b) Analyze the Indian Constitutional Triad (Articles 21, 48A, and 51A(g)) and compare the Chipko Movement with the Narmada Bachao Andolan. [5 Marks]
}{10}{CO5, CO6}{L4 Analyze}

\newpage
\subsection*{Solutions and Marking Rubrics --- End-Term Set B}

\begin{solutionbox}[Solution to Descriptive Question 3(b) --- Biodiversity Indices]
\textbf{1. Relative Abundances ($p_i = n_i / N$ with $N = 100$):}\\
$p_1 = 0.40, \; p_2 = 0.30, \; p_3 = 0.20, \; p_4 = 0.10$.\\
\textbf{2. Shannon-Wiener Index $H'$ (2.5 Marks):}\\
$p_1 \ln p_1 = 0.40 \times (-0.9163) = -0.3665$\\
$p_2 \ln p_2 = 0.30 \times (-1.2040) = -0.3612$\\
$p_3 \ln p_3 = 0.20 \times (-1.6094) = -0.3219$\\
$p_4 \ln p_4 = 0.10 \times (-2.3026) = -0.2303$\\
$H' = -\sum p_i \ln p_i = -(-1.2799) = 1.280\text{ nats}$.\\
\textbf{3. Simpson's Index of Diversity (2.5 Marks):}\\
$D = \sum p_i^2 = (0.40)^2 + (0.30)^2 + (0.20)^2 + (0.10)^2 = 0.16 + 0.09 + 0.04 + 0.01 = 0.30$.\\
$\text{Simpson Diversity } = 1 - D = 1 - 0.30 = 0.70$.
\end{solutionbox}

\begin{table}[h]
\centering\small
\begin{tabular}{|c|c||c|c||c|c||c|c||c|c|}
\hline
\textbf{Q} & \textbf{Key} & \textbf{Q} & \textbf{Key} & \textbf{Q} & \textbf{Key} & \textbf{Q} & \textbf{Key} & \textbf{Q} & \textbf{Key} \\ \hline

1 & C & 7 & C & 13 & B & 19 & A & 25 & A \\ \hline
2 & A & 8 & D & 14 & D & 20 & C & 26 & B \\ \hline
3 & B & 9 & A & 15 & A & 21 & C & 27 & C \\ \hline
4 & D & 10 & C & 16 & B & 22 & A & 28 & D \\ \hline
5 & A & 11 & C & 17 & C & 23 & B & 29 & A \\ \hline
6 & B & 12 & A & 18 & D & 24 & D & 30 & C \\ \hline
\end{tabular}
\caption{Official Answer Key for Part A Objective MCQs --- End-Term Set B}
\end{table}

\newpage
